\pagebreak

\chapter{Target Statements}\label{chap:target}

    Target statements are copied to a target code structure during the
    immediate execution. The immediate code selects which statements are
    copied and may duplicate target code with variations introduced by
    immediate expressions.
    
    Note that \\
    {\bf value mode assignment}, \\
    {\bf priority mode assignment} , \\
    {\bf clock mode assignment} and \\
    {\bf unconditional output mode assignment} \\
    do not generate in-line target code and so
    have been included under immediate assignment above
    \autorefph{chap:immed}.

    \section{expression availability and operators}

        \index{queue!availability}
        An expression is available when the contained queue variables are
        available. This means that the conditional operator {\bf ?:} cannot be
        used to select an available subexpression where the alternative is not
        available. For example in the following case -
        \begin{lstlisting}[gobble=8, language=threepl]
               q1 <<- l ? q2 : q3;
        \end{lstlisting}
        where {\bf q1} to {\bf q3} are queues, all three queues must be available
        for the statement to execute even though only one of queues {\bf q2} and
        {\bf q3} will be read. Similarly if this statement is enclosed in a
        {\bf sync} control structure the sync body will not be executed
        until all three queues are avaliable regardless of conditional selection
        by variable {\bf l}.

    \section{selectvalue mode assignment}\label{sec:target_selvalass}

        \index{assignment!selectvalue mode}
        A selectvalue assignment statement has the form -
        \begin{lstlisting}[gobble=8, language=threepl]
               variable |- expression;
        \end{lstlisting}
        
        Whole structs or arrays can be assigned by
        \begin{lstlisting}[gobble=8, language=threepl]
               variable |- variable;
        \end{lstlisting}

        \subsubsection{immediate execution}
        
            On immediate execution this generates one in-line target
            statement.
        
        \subsubsection{target execution}
        
            Selectvalue assignment is assignment to a selectvalue variable.

            RHS expressions may be immediate or target mode.

            A selectvalue variable represents a bus, each assignment being a
            drive onto the bus\footnote{CLB logic is used by default but where
            the target device has internal 3-state buffers these can be used to
            implement selectvalue variables. The attribute "threestate" is used
            to force the use of 3-state buffers.}. This is a target assignment
            in that the assignment is only valid for one clock cycle - the bus
            is driven for one cycle when this statement executes. If the RHS
            contains one or more queue reads execution may be delayed waiting for
            queue availability \index{queue!availability}. The assignment can be
            held if required by repeated execution.
            
            For compound types assignments can be made to individual fields
            or array members. On any clock cycle when a component of the
            variable is not assigned the output for that component will be
            the default value specified when the selectvalue variable was
            created.
        
            A selectvalue variable has a single clock domain for reading and
            writing. When a selectvalue variable is created using the inbuilt
            procedure {\bf selectvalue()} \autorefph{sec:ipselectvalue} it has
            no associated clock domain. The clock domain is established by
            \blist
            \item   the first assignment to the variable (or component of
                    the variable if the type is compound)
            \item   the first use of the value
                    of the variable or one of its components
            \blend
            The clock domain of the variable will be set to the current clock
            domain.
            
            If a value to be assigned to a selectvalue variable has a different
            clock domain it must be converted by one of the inbuilt functions
            {\bf sample()} \autorefph{sec:ifsample} or {\bf accept()}
            \autorefph{sec:ifaccept}. Note that evaluation of the variable in
            any context, for example in inbuilt or library modules, procedures
            or functions, will establish a clock domain if the variable does not
            yet have a clock domain. This will always be the current clock
            domain.

            If the RHS expression contains queue variables their output (read)
            clock domain must be the current clock domain (if this is the first
            read or examine for a queue its output domain will be established by
            the assignment).

            Whether the assignment is of a primitive or a compound type the
            widths of individual words may differ as long as the primitive types
            match. Bits and unsigned integers will be padded or truncated as
            needed. Signed integers will be truncated or sign extended. Signed
            integers may not be assigned to unsigned integers. However, unsigned
            integers may be assigned to signed integers, since size adjustments
            are made by truncation or zero padding the sign of the result may be
            incorrect unless the LHS is at least one bit larger than the RHS.





    \section{static mode assignment}\label{sec:stat_ass}

        \index{assignment!static mode}
        Assignment to a static variable has the forms
        \begin{lstlisting}[gobble=8, language=threepl]
               variable <- expression;
               variable +<- expression;
               variable -<- expression;
               variable *<- expression;
               variable++;
               variable--;
        \end{lstlisting}
        Where a variable is a structure or an array, fields or members can be
        assigned by
        \begin{lstlisting}[gobble=8, language=threepl]
               variable.field <- expression;
               variable[subscript] <- expression;
               variable.field +<- expression;
               variable[subscript] +<- expression;
               variable.field -<- expression;
               variable[subscript] -<- expression;
               variable.field *<- expression;
               variable[subscript] *<- expression;
               variable.field++;
               variable.field--;
               variable[subscript]++;
               variable[subscript]--;
        \end{lstlisting}
        or any depth of nesting of fields and members can be used. Whole
        structures or arrays can be assigned by
        \begin{lstlisting}[gobble=8, language=threepl]
               variable <- variable;
        \end{lstlisting}
        
        The increment and decrement statement forms can only be applied to types
        "int:" and "uint:". The \verb'+<-', \verb'-<-' and \verb'*<-' forms can
        only be applied to types "int:", "uint:", "fixed:" and "ufixed:". Type
        conversions in assignment are allowed in many cases and these are
        described in \autorefph{sec:assignments}.
        

        \subsubsection{immediate execution}
        
            On immediate execution this generates one in-line target statement.
        
        \subsubsection{target execution}        

            RHS expressions and variables may contain queues or queue fields
            which are read or examined.

            Regardless of any context or enclosing constructs which control
            execution, a static assignment can only execute if the RHS
            expression is available.

            If the LHS variable is a struct or an array several assignment
            statements may be needed. These assignment statements may be
            grouped to execute simultaneously.

            Where assignment to a static variable which is a struct or array
            is incomplete, the unassigned fields or array members will be
            unchanged.

            Whether the assignment is of a primitive or a compound type the widths
            of individual words may differ as long as the primitive types match.
            Bits and unsigned integers will be padded or truncated as needed. Signed
            integers will be truncated or sign extended. Signed integers may not be
            assigned to unsigned integers. However, unsigned integers may be
            assigned to signed integers, since size adjustments are made by
            truncation or zero padding the sign of the result may be incorrect
            unless the LHS is at least one bit larger than the RHS.

            In any one execution cycle the same array member or struct field
            may not be assigned to by more than one assignment statement. Such
            conflicts cannot be identified in general by the compiler. If an
            assignment conflict occurs the resulting value is undefined.
            \footnote{At the moment the compiler does not detect any
            assignment conflicts. It may in a future version. Assignment
            conflicts could also be detected at execution time and cause
            exceptions to be thrown, but that has not been implemented yet.}
            
            Optimisation note: if an immediate value, or a variable
            which is value mode and is constant, is assigned to a static
            variable or one of its components, the constant is checked
            against the initial value for that variable or component. If
            they are the same the variable or component is reset to its
            initial value rather than being assigned. This has
            benefits in both gate usage and speed -
            \blist
            \item   with one fewer assignment the data input logic is simplified
                    using fewer gates
            \item   simpler data input logic means that it will usually be faster
            \item   reset logic is inherently simpler, using fewer gates than
                    data input logic
            \item   being simpler reset logic is inherently faster
            \blend
            
            A static variable, or one or more components of a variable, can be
            explicitly reset by using the inbuilt procedure
            {\bf reset()} \autorefph{sec:ipreset}
                              
            A static variable has a single clock domain for reading and
            writing. When a static variable is created using the inbuilt
            procedure {\bf static()} \autorefph{sec:ipstatic} it has no
            associated clock domain. The clock domain is established by -
            \blist
            \item   the first assignment to the variable (or component of
                    the variable if the type is compound)
            \item   the first use of the value
                    of the variable or one of its components
            \blend
            The clock domain of the variable will be set to the current clock
            domain.

            {\bf Note that evaluation of the variable in any context, for example
            in inbuilt or library modules, procedures or functions, will establish
            a clock domain if the variable does not yet have a clock domain.}
            This will always be the current clock domain (unless the code has
            been explicitly written to establish the clock domain of its argument
            to something other than the current clock domain, in which case
            this will presumably be well documented).

            If the {\bf reset()} inbuilt procedure is used on a static variable -
            \blist
            \item   if the static variable has a clock domain established, it
                    must be the same as the current clock domain
            \item   if the static variable does not have a clock domain established,
                    it will be given the current clock domain
            \blend
            
            An assignment executes on one clock cycle unless the RHS
            expression contains queue reads on queues which are not available.



    \section{queue mode assignment}\label{sec:queue_ass}

        \index{assignment!queue mode}
        Assignment to a queue variable is of the form
        \begin{lstlisting}[gobble=8, language=threepl]
               variable <<- expression;
        \end{lstlisting}
        Where a variable is a structure or an array, fields or
        members can be assigned by
        \begin{lstlisting}[gobble=8, language=threepl]
               variable.field <<- expression;
               variable[subscript] <<- expression;
        \end{lstlisting}
        or any depth of nesting of fields and members can be used. Whole
        structures or arrays can be assigned by
        \begin{lstlisting}[gobble=8, language=threepl]
               variable <<- variable;
        \end{lstlisting}

        \subsubsection{immediate execution}
        
            On immediate execution this generates one in-line target statement.
        
        \subsubsection{target execution}
        
            The {\bf \verb'<<-'} assignment writes to a queue. The datum
            becomes available on the queue output.
            
            An assignment executes on one clock cycle unless the RHS
            expression contains queues which are not available (are empty) or
            the LHS queue is not available (is full).
        
            A queue variable has a clock domain for reading and another clock
            domain for writing. When a queue variable is created using the
            inbuilt procedure {\bf queue()} \autorefph{sec:ipqueue} it has no
            associated clock domains.

            The clock domain for writing is established by -
            \blist
            \item   the first write (assignment) to the queue variable (or
                    component of the variable if the type is compound)
            \item   the use of the write availability operator \verb'>' on the
                    queue
            \item   a call to {\bf queuewritestatus()} on the queue
            \item   a call to {\bf reset()} on the queue
            \blend
            If the write clock domain is not yet established in any of these
            situations then it will be set to the current clock domain,
            otherwise it must be the same as the current clock domain.

            The clock domain for reading is established by -
            \blist
            \item   the first read or examine of the variable (or component of
                    the variable if the type is compound)
            \item   the use of the read availability operator \verb'<' on the
                    queue
            \item   a call to {\bf queuereadstatus()} on the queue
            \blend
            If the read clock domain is not yet established in any of these
            situations then it will be set to the current clock domain,
            otherwise it must be the same as the current clock domain.

            The read and write clock domains will usually be the same, in which
            case the queue variable is called {\bf synchronous}.

            When a queue variable is written (assigned), If it has no write
            clock domain the current clock domain will be established as the
            write clock domain of the variable. If it has a write clock domain
            it must be the same as the current clock domain.

            When a queue variable is written (assigned), the clock domain of the
            assigned value must be the same as the current clock domain or
            null.

            If the {\bf reset()} inbuilt procedure is used on a queue variable -
            \blist
            \item   if the queue variable has a write clock domain established,
                    it must be the same as the current clock domain
            \item   if the queue variable does not have a write clock domain
                    established, it will be given the current clock domain
            \blend

            When a queue variable read or examined the resulting value will have
            the read clock domain of the queue variable.

            RHS expressions and variables may contain queues or queue fields,
            which may be read or examined.

            Regardless of any context or enclosing constructs which control
            execution, a queue assignment can only execute if the LHS queue is
            available and the RHS expression is available.

            When a queue of compound data type is accessed the entire compound
            datum is read or written. On a read any unused struct fields or
            array members are ignored. On a write any unassigned struct fields
            or array members will be clear (0, false etc.)\footnote{For an
            enumerated type the value with an ordinal of 0 will be assigned. If
            the enumerated type has been created using enumv(), rather than
            enum(), such that there is no value for ordinal 0 then the assigned
            value will be invalid!}.

            Compound types or subtypes can be assigned in a single assignment
            statement where the RHS can be expressed as a single variable or
            part of a variable. For example 
            \begin{lstlisting}[gobble=12, language=threepl]
                   queue(p, "[4][2]uint:8");
                   static(s, "[4][4]uint:10");

                   p <<- s[][1..2];
            \end{lstlisting}
            will assign a subarray of static variable {\em s} to queue {\em p}
            and will truncate the 10 bit unsigned integers to 8 bit.

            Members of a compound data type can be accessed in separate
            assignment statements as long as those statements are executed
            simultaneously so that a single read or write occurs. Simultaneity
            can be enforced by enclosing these statements in a {\em sync}
            block. For example in the following
            \begin{lstlisting}[gobble=12, language=threepl]
                   struct(s,
                       f1,     "log",
                       f2,     "uint:8"
                   );
                   queue(q1, s);
                   queue(q2, s);

                   q1.f1 <<- q2.f1;
                   q1.f2 <<- q2.f2 + 3;
            \end{lstlisting}
            the two statements would execute simultaneously since the queue
            dependencies are the same (q1 and q2). Where different queues are
            used we need to synchronise the statements, e.g.
            \begin{lstlisting}[gobble=12, language=threepl]
                   struct(s,
                       f1,     "log",
                       f2,     "uint:8"
                   );
                   queue(q1, s);
                   queue(q2, s);
                   queue(q3, s);

                   sync {
                       q1.f1 <<- q2.f1;
                       q1.f2 <<- q3.f2;
                   }
            \end{lstlisting}
            
            In any one execution cycle the same array member or struct field may
            not be assigned to by more than one assignment statement. Such
            conflicts cannot be identified in general by the compiler. If an
            assignment conflict occurs the resulting value is undefined
            (actually the data are ORed - also see the footnote in the
            discussion of assignment conflicts for static assignments in the
            previous section.) It is usual to rely on the logic of the code and
            its timing to ensure mutual exclusion but in the case of data
            resulting from multiple asynchronous external sources this may not
            be possible. In this case the waitpri() control structure, section
            \autorefph{sec:waitpri}, can be used to arbitrate.

            \index{null!queue}
            \index{queue!null}
            A queue may have data type "null", e.g.
            \begin{lstlisting}[gobble=12, language=threepl]
                   queue(trigger, "null");
            \end{lstlisting}
            in which case it will have no
            associated data. Such a queue can be used to signal an event.
            A queue whose data type is {\em null} can be written by
            \begin{lstlisting}[gobble=12, language=threepl]
                   variable <<- null;
            \end{lstlisting}
            and can be read by
            \begin{lstlisting}[gobble=12, language=threepl]
                   null <<- variable;
            \end{lstlisting}
            
            A queue of type null can be assigned to another queue of type null.
            
            A queue of type null cannot return a value on reading, however
            it can be used as if it were type "log" by casting it to
            "log", in which case the value returned will always be {\bf true}.
            
            In the case of the target conditional statements {\bf when} and
            {\bf sync when} the conditional expression can be a single queue
            mode variable of type "null" which will be treated as though of
            type "log" whose value will always evaluate as TRUE.            
            
            Thus a wait on a null queue can be explicit as in -
            \begin{lstlisting}[gobble=12, language=threepl]
                   when (trigger) sync {
                       .
                       .
                       .
                   }
            \end{lstlisting}
            Alternatively an assignment to null can be used. In the following
            example it is enclosed in a sync block so will wait until the
            queue is available and will then explicitly acknowledge (read)
            the queue by assignment to null -
            \begin{lstlisting}[gobble=12, language=threepl]
                   sync {
                       null <<- trigger;
                       .
                       .
                       .
                   }
            \end{lstlisting}
            In each case null queue {\bf trigger} will be waited upon and
            then read.
            
            In other contexts a null queue may be cast to type "log" where
            again its value will always be true.

            To acknowledge queues whose values have been extracted by the
            examine operator {\bf ?} assignment to null can be used (read and
            ignore the result) -
            \begin{lstlisting}[gobble=12, language=threepl]
                   null <<- expression;
            \end{lstlisting}
            where if the expression contains any queues they will be
            acknowledged without their values being read. It would be usual
            for the RHS to be either a queue variable or a value variable
            representing an expression. In the latter case any queue reads
            contained in the expression would occur and be acknowledged.
            An expression, rather than a variable, occurring on the RHS
            would generate unused code (though this would be discarded by
            the code generator optimiser) and hence is pointless, though
            not erroneous.

            Whether the assignment is of a primitive or a compound type the
            widths of individual words may differ as long as the primitive
            types match. Bits and unsigned integers will be padded or
            truncated as needed. Signed integers will be truncated or sign
            extended. Signed integers may not be assigned to unsigned
            integers. Unsigned integers may be assigned to signed integers,
            but since size adjustments are made by truncation or zero
            padding the sign of the result may be incorrect unless the
            LHS is at least one bit larger than the RHS.







    \section{output mode conditional assignment}\label{sec:target_outputtass}

        \index{assignment!selectvalue mode}
        Conditional assignment of output mode variables is used where 3-state
        output is required.
        
        Unconditional assignment to an output mode variable was described
        previously (\autorefph{sec:outputiass}).

        Conditional output mode assignment can have the form -
        \begin{lstlisting}[gobble=8, language=threepl]
        when (condition)
               variable |- expression;
        \end{lstlisting}
        where the {\bf when} is at the top level, i.e. not nested within any other
        control structure except a module. It is preferable that the condition be
        simply a static variable of type "log" with the {\bf direct} and {\bf iob}
        attributes set to true since in that case the static variable flip-flops
        can be placed in the IO blocks for better timing.
        
        A conditional assignment may also be executed within any normal execution
        thread, e.g.
        \begin{lstlisting}[gobble=8, language=threepl]
        seq {
            .
            .
            variable |- expression;
            .
            .
        }
        \end{lstlisting}
        
        In each of the above the pad output will be in high impedance state except
        for the clock cycle during which the conditional assignment statement
        executes.
                
        The LHS variable may have a compound type and the assignment may reference
        all or part of that variable. Assignments to components of a compound type
        may be conditional, unconditional or a mixture of both.
        
        Conditional or unconditional assignment to a component of a variable can
        only be done once.
        
        An output mode variable must have at least a {\bf loc} attribute,
        providing a location for each pad.

        \subsubsection{immediate execution}
        
            On immediate execution this generates one in-line target
            statement.
        
        \subsubsection{target execution}

            RHS expressions may be immediate or target mode. If the value to be
            assigned has a different clock domain it must be converted by one of
            the inbuilt functions {\bf sample()} \autorefph{sec:ifsample} or {\bf
            accept()} \autorefph{sec:ifaccept}.
            
            A RHS expression cannot contain a queue read (but can contain a queue
            examine or a queue availability operator).

            An output variable represents a connection off-chip via output buffers
            and pads. The buffers may be 2-state or 3-state.
            
            In the Xilinx implementation if the source of the output is a static
            variable or an examined synchronous queue variable then the associated
            flip-flops can be located in the IOBs. This can be done by a suitable
            flag to the Xilinx tools but is best done by explicitly setting the IOB
            attribute on the static variable since the 3PL netlist optimiser is
            then aware of this. If the flip-flops go to other destinations as well
            as the output buffers then they will be replicated, the replicants
            being connected in parallel and located in IOBs while the originals
            stay in CLBs. Similarly if the expression controlling conditional
            output is a static variable or an examined synchronous queue variable
            then the associated flip-flop will be replicated for each output bit,
            the replicants being connected in parallel and located in IOBs.
            Placement of output flip-flops in IOBs removes the necessity to apply
            timing constraints to these outputs. In light of the above, in the
            following code -
            \begin{lstlisting}[gobble=12, language=threepl]
                   when (l)
                       o |- v;
            \end{lstlisting}
            {\bf l} and {\bf v} should be static variables if possible and {\bf o}
            should have attribute {\bf direct} set. As described above in that case
            the flip-flops for the variable will be located in the IOBs and the
            propagation time from flip-flop to pad will be the minimum achievable
            (controlled also by slew rate and drive current attributes). If those
            conditions are not met then a timing constraint from the clock domain
            to the output pads will usually need to be applied.

            Whether the assignment is of a primitive or a compound type the widths
            of individual words may differ as long as the primitive types match.
            Bits and unsigned integers will be padded or truncated as needed.
            Signed integers will be truncated or sign extended. Signed integers may
            not be assigned to unsigned integers. Unsigned integers may be assigned
            to signed integers, but since size adjustments are made by truncation
            or zero padding the sign of the result may be incorrect unless the LHS
            is at least one bit larger than the RHS.
        
            An output variable has a write clock domain but no read clock domain.
            When an output variable is created using the inbuilt procedure {\bf
            output()} \autorefph{sec:ipoutput} it has no associated clock domain. Its
            clock domain is established by the first assignment to the variable (or
            component of the variable if the type is compound) and is the current
            clock domain. The RHS exression must also have the current clock domain
            as its output (read) clock domain.










    \section{the {\bf seq}, {\bf par} and {\bf sync} control structures}\label{sec:seqparsync}

        \index{target!seq}
        \index{target!par}
        \index{target!sync}
        These are of the form -
        \begin{lstlisting}[gobble=8, language=threepl]
           seq {
               body
           }
        \end{lstlisting}
        or
        \begin{lstlisting}[gobble=8, language=threepl]
           par {
               body
           }
        \end{lstlisting}
        or
        \begin{lstlisting}[gobble=8, language=threepl]
           sync {
               body
           }
        \end{lstlisting}
        The body is executed sequentially during immediate execution but any
        target statements encountered in the body generate code for the
        target which will execute them sequentially (seq) or in parallel (par
        and sync).

        Target code in a {\bf seq \{\dots\}} block will be executed
        sequentially in the target. It may contain any blocks or statements.
        Each component of the body will execute when the previous one has
        completed. Any components may wait on queues. In the target the block
        will complete when the last component has finished execution.

        Target code in a {\bf par \{\dots\}} block will be executed in
        parallel. It may contain any blocks or statements. All components
        will execute once. Any component may wait on queues. The block will
        complete when all components have executed.

        \index{queue!availability}
        Target code in a synchronised block {\bf sync \{\dots\}} will be
        executed in parallel when all queues read or written (but not examined
        or availability tested) anywhere in the block are available. This
        includes any queues read or written in procedures or functions called.
        Note that if a queue read is conditional this is not taken into account
        and the sync will still wait on read availability.
        
        Note that the {\em sync} only controls initial entry into the parallel
        block and not any subsequent behaviour. For that reason it would be
        unusual, though not erroneous, to include other blocks or loops within a
        {\em sync} block.

        Using a {\em sync} block can improve code by taking the queue
        availability test away from individual statements in the block and
        using a single test instead. In many cases the compiler recognises
        that code is within a {\em sync} block and omits the queue
        availability test. This is done with care since later queue reads in
        say a {\em seq} block within a {\em sync} block need to check queue
        availability since another read may have been done since the {\em
        sync} availability test.








    \section{the {\bf when} control structure}\label{sec:when_struct}
    
        \index{statement!when}
        The {\bf when} has the formats -
        \begin{lstlisting}[gobble=8, language=threepl]
               when (expression)
                   body
        \end{lstlisting}
        and
        \begin{lstlisting}[gobble=8, language=threepl]
               when (expression)
                   true_body
               else
                   false_body
        \end{lstlisting}

        The expression must evaluate to a logical result (see exception below)
        and is usually target mode, but it may be immediate mode. If the test
        expression is a constant value mode variable, or is immediate mode, then
        the test will be omitted in the target code and the true or false block
        will be created as appropriate. If the test expression is not a constant
        value mode variable then it must have a single clock domain which is the
        same as the current clock domain.
        
        The test expression may be a queue of type "null" in which case it will
        be treated as though of type "log" whose value is true. This is treated
        as a special case and hence a null queue cannot occur within a logical
        expression without casting it to type "log".
        
        The code bodies are a single statement, which may be a {\em seq},
        {\em par} or {\em sync} block. A code body may be a plain block
        {\em \{\dots\}} as long as it does not contain more than one target
        statement.
        
        Either of the above forms of {\bf when} may be followed by an
        exception termination
        \begin{lstlisting}[gobble=8, language=threepl]
               unless (exception_expression) {}
        \end{lstlisting}
        or
        \begin{lstlisting}[gobble=8, language=threepl]
               unless (exception_expression) {
                   exception_body
               }
        \end{lstlisting}
        
        Exception termination constructs are allowed on several control structures
        - see \autorefph{sec:excep_term}.
        
        \subsection{immediate execution}
        
            On immediate execution this generates an in-line target statement.
            Immediate sequential execution is carried out on the code body or
            bodies of the {\em when}.
        
        \subsection{target execution}

            In the first form the target body will be executed if the
            expression is true. In the second form one of the target code
            bodies will be executed as selected. Following either of the above
            the {\bf when} will complete execution.

            If the test expression contains queue variables the test expression
            may wait on queue availability \index{queue!availability}. If the test
            expression reads a queue it will remove a datum. If a statement
            within an enclosed code body also reads the same queue, care must be
            taken to ensure that these reads occur simultaneously, otherwise two
            data will be read rather than one (see examples below). In some
            cases it may be best to use an examine rather than a read on a queue
            variable within a test expression for a control structure.
            
            The {\em when} test adds no execution cycles so the {\em when}
            executes in the number of cycles taken by the code body executed,
            with the exception that a missing code body will execute in one
            clock cycle (for example if there is no {\em else} body and the
            test is false then the {\em when} will take one clock cycle).
            
            The following is an example of the correct use of queues -
            \begin{lstlisting}[gobble=12, language=threepl]
                   when (p == 3)
                       s <- p;
            \end{lstlisting}
            The {\bf when} test will wait on the availability of queue {\bf p}.
            Once {\bf p} is available its value will be tested. If the value
            of the queue variable is 3, the assignment to the static variable
            {\bf s} will be performed. If the test fails queue {\bf p} will
            nevertheless have been read during evaluation of the test 
            expression. If the test succeeds, the queue read in the test
            expression and the queue read in the enclosed assignment statement
            will take place in the same clock cycle and hence only a single read
            of one datum will take place.
            
            The following example is incorrect -
            \begin{lstlisting}[gobble=12, language=threepl]
                   when (q1 == 3)
                       s <- q1 + q2;
            \end{lstlisting}
            The {\bf when} test will wait on the availability of queue {\bf q1}.
            Once {\bf q1} is available its value will be tested and if
            successful the assignment will be executed. The assignment statement
            may now wait on the availability of queue {\bf q2}, but the read of
            queue {\bf q1} in the test expression will have gone ahead
            regardless. When queue {\bf q2} becomes available and queue {\bf q1}
            again becomes available {\bf q1} will be read again. This example is
            easily fixed by enclosing the {\bf when} in a sync block -
            \begin{lstlisting}[gobble=12, language=threepl]
                   sync {
                       when (q1 == 3)
                           s <- q1 + q2;
                   }
            \end{lstlisting}
            which will wait to enter the {\bf when} until all enclosed queues
            are available, at which time the {\bf when} will execute both its
            test evaluation and conditionally the enclosed statement in one
            clock cycle. Alternatively it could have been written as -
            \begin{lstlisting}[gobble=12, language=threepl]
                   when (<q1 && ?q1 == 3)
                       s <- q1 + q2;
            \end{lstlisting}
            which would do an examine of {\bf q1}, rather than a read, in the
            test expression (having first verified availability of {\bf q1}).
            This is less elegant than the previous solution and is less
            efficient.
            
            The compiler recognises cases where a {\bf when()} contains
            queue reads or writes in one of the code bodies and also reads
            some but not all of these queues in the test expression and
            issues a warning.

            Note that {\bf when} statements pass execution control
            through the test expression within the same cycle that the body
            executes. There are two consequences of this that should be noted.
            
            Firstly the signal delay through the test expression may reduce
            the timing margins for the first statement (seq block) or statements
            (par block) in the code bodies. In some cases it may for example
            be possible to change the order of statements in a seq block so that
            a complex statement is not the first.
            
            Secondly nested {\bf when} statements concatenate the control delay
            through each of their test expressions. For this reason the depth to
            which it is prudent to nest {\bf when} statements depends on the clock
            frequency. At high clock frequencies it is preferable to use multiple
            parallel {\bf when} statements, for example -
            \begin{lstlisting}[gobble=12, language=threepl]
                   when (l1)
                       when (l2)
                           statement1
                       else
                           statement2
            \end{lstlisting}
            may be better written as -
            \begin{lstlisting}[gobble=12, language=threepl]
                   when (l1 && l2)
                       statement1
                   when (l1 && !l2)
                       statement2
            \end{lstlisting}
            assuming that these are in a parallel execution context (and
            assuming that the clock frequency is too high for the nested
            example). Note that if a logical expression is complex, e.g. -
            \begin{lstlisting}[gobble=12, language=threepl]
                   when (i < 4567)
                       when (l2)
                           statement1
                       else
                           statement2
            \end{lstlisting}
            its duplication can be avoided by -
            \begin{lstlisting}[gobble=12, language=threepl]
                   value(l, "log", i < 4567);
                   when (l && l2)
                       statement1
                   when (l && !l2)
                       statement2
            \end{lstlisting}
                    
            If an exception termination construct is appended, execution within
            the {\bf when} can be prematurely terminated
            - see \autorefph{sec:excep_term}.







    \section{the {\bf sync when} control structure}\label{sec:sync_when_struct}
    
        \index{statement!sync when}
        The {\bf sync when} has the format -
        \begin{lstlisting}[gobble=8, language=threepl]
               sync when (expression) {
                   body
               }
        \end{lstlisting}

        The expression must evaluate to a logical result (see exception below)
        and is usually target mode, but it may be immediate mode. If the test
        expression is a constant value mode variable, or is immediate mode, then
        the test will be omitted in the target code and the code block will be
        created or not as appropriate. If the test expression is not a constant
        value mode variable then it must have a single clock domain which is the
        same as the current clock domain. The test expression cannot contain
        unbuffered queue reads.

        The test expression may be a queue of type "null" in which case it will
        be treated as though of type "log" whose value is true. This is treated
        as a special case and hence in general a null queue cannot occur within
        a logical expression without casting it to type "log".

        The code body in braces is a parallel block as in the case of {\bf sync}
        and {\bf sync while}. The code block cannot contain unbuffered queue
        reads or writes.
        
        The {\bf sync when} may be followed by an
        exception termination
        \begin{lstlisting}[gobble=8, language=threepl]
               unless (exception_expression) {}
        \end{lstlisting}
        or
        \begin{lstlisting}[gobble=8, language=threepl]
               unless (exception_expression) {
                   exception_body
               }
        \end{lstlisting}
        
        Exception termination constructs are allowed on several control structures
        - see \autorefph{sec:excep_term}.
        
        \subsection{immediate execution}
        
            On immediate execution this generates an in-line target statement.
            Immediate sequential execution is carried out on the code body of
            the {\em sync when}.
        
        \subsection{target execution}

            The {\bf sync when} will execute the enclosed code block if the
            test expression is true and all queue variables in the test
            expression and code block are available. If not, execution of the
            statement terminates. This differs from a {\bf when} inside a 
            {\bf sync} where the {\bf sync} will wait on queue availability.
            If the {\bf sync when} fails execution of the thread continues
            with no delay, i.e. the statement takes no clock cycles to execute.
            
            It is important to note that the {\bf sync} takes no account
            of conditional access to queues within its block, i.e. -
            \blist
            \item   queue reads or writes within a {\bf when} statement
            \item   queue reads within operands to the ternary conditional
                    operator \verb'?:'
            \blend
            If queue reads or writes are located within the block the {\bf sync}
            will test for availability on all of these before execution of the
            target code in the block, even though some of these queue accesses
            may not eventuate due to conditional execution.
            
            For example in the following case -
            \begin{lstlisting}[gobble=12, language=threepl]
                queue({q1, q2", "log");
                sync when (q1) {
                    seq {
                        q1 <<- q2;
                        q2 <<- q3;
                    }
                }
            \end{lstlisting}
            the statement will only succeed if queue q2 contains a single datum
            since it has the default depth of 2. If q2 is empty \verb'<q2' will be
            false and if q2 is full \verb'>q2' will be false. No account is taken
            of the fact that q2 is read before it is written.
            
            The target code body will be executed if the expression is true and all
            queues read in the test expression or read or written anywhere in the
            code body are available. Statements in the code body are executed in
            parallel. If any queues to be read in the test expression, or any queues
            in the code body, are not available the {\bf sync when} will not
            evaluate the test expression and hence will not execute the code body
            and execution of the {\bf sync when} will finish. Similarly if
            evaluation of the test expression yields a false result the code body
            will not be executed.

            If the test expression contains queue variables the test expression
            will not wait on their availability \index{queue!availability}. No
            queues in the test expression will be read unless all queues to be
            read in the test expression are available and all queues to be read
            and written in the code block are available.

            The {\em sync when} test adds no execution cycles so the {\em sync
            when} executes in the number of cycles taken by the code body.

            Note that {\bf sync when} statements pass execution control
            through the test expression within the same cycle that the body
            executes. The signal delay through the test expression may reduce
            the timing margins for the first statement (seq block) or statements
            (par block) in the code bodies. In some cases it may for example
            be possible to change the order of statements in a seq block so that
            a complex statement is not the first.
                    
            If an exception termination construct is appended, execution within
            the {\bf sync when} can be prematurely terminated
            - see \autorefph{sec:excep_term}.
            
            The {\bf sync when} is intended to be used where queue waits are to be
            avoided. For example to control queue assignment one might write -
            \begin{lstlisting}[gobble=12, language=threepl]
                   when (go)
                       q2 <<- q1;
            \end{lstlisting}
            The disadvantage of this is that if {\bf go} is true but {\bf q1} is not
            available for reading, or {\bf q2} is not available for writing, then
            the assignment statement will hang waiting for queue availability.
            If we attempt to avoid entering the {\bf when} by enclosing it in
            a {\bf sync} -
            \begin{lstlisting}[gobble=12, language=threepl]
                   sync {
                       when (go)
                           q2 <<- q1;
                   }
            \end{lstlisting}
            the {\bf sync} itself will hang waiting for availability. We could
            change the first attempt to -
            \begin{lstlisting}[gobble=12, language=threepl]
                   when (>q2 && <q1 && go)
                       q2 <<- q1;
            \end{lstlisting}
            which achieves what we want. The {\bf when} will not execute the assignment
            unless the two queues are available and {\bf go} is true. If these
            conditions are not met the {\bf when} finishes and passes execution on.

            Using a {\bf sync when} we could instead write -
            \begin{lstlisting}[gobble=12, language=threepl]
                   sync when (go) {
                       q2 <<- q1;
                   }
            \end{lstlisting}
            The {\bf sync when} achieves the same result as the final example above
            while avoiding the necessity of finding all queue reads and writes in
            the code block and explicitly testing for their availability.
            
            If in the above example the variable {\bf go} was a queue then the last
            two examples are no longer equivalent. A queue read encountered in the
            test expression of a {\bf when} will cause the {\bf when} to hang
            waiting for the queue to become available before it can evaluate its
            test expression. The {\bf sync when} will abandon any attempt to evaluate the
            test expression in that case since it can see that the queue is not
            available. It could notionally be coded as -
            \begin{lstlisting}[gobble=12, language=threepl]
                   when (<go)
                       when (go && >q2 && <q1)
                           q2 <<- q1;
            \end{lstlisting}
            While this is functionally equivalent to the use of {\bf sync when} as
            listed above, it generates logic which is much less efficient, both
            in path delays and resource usage.





    \section{the {\bf branch} control structure}\label{sec:branch_struct}
    
        \index{statement!branch}
        The {\bf branch} has the format -
        \begin{lstlisting}[gobble=8, language=threepl]
               branch (expression) {
               case cexpr1:
                   body1
               case cexpr2:
                   body2
               default:
                   body3
               }
        \end{lstlisting}
        
        This branches to one of the case code bodies as a result of testing
        the value of the expression, which must be a target mode. The case
        expressions must be immediate mode. A case body target code does not flow
        on to the next case body target code.
        
        A case may have a comma separated list of immediate expressions
        within braces, e.g.
        \begin{lstlisting}[gobble=8, language=threepl]
               branch (expression) {
               case {cexpr1, cexpr2, cexpr3}:
                   body1
               case cexpr4:
                   body2
               }
        \end{lstlisting}
        in which case if the branch expression matches any of the case
        expression values the code body will be executed.

        Case values are immediate expressions which evaluate to constants for
        the target code. The resulting case values must be unique. Primitive
        types allowed are "bits", "uint", "int" and "enum".

        The branch statement terminates when the selected code body terminates.

        The default case is optional. If no cases are matched then no target
        statements are executed and the branch statement terminates.

        The branch expression and the case expressions must be compatible types.
        The test expression must have a single clock domain which is the same as
        the current clock domain.

        The code bodies are single statements, which may of course be {\em seq},
        {\em par} or {\em sync} blocks.
        
        The {\bf branch} may be followed by an exception termination
        \begin{lstlisting}[gobble=8, language=threepl]
               unless (exception_expression) {}
        \end{lstlisting}
        or
        \begin{lstlisting}[gobble=8, language=threepl]
               unless (exception_expression) {
                   exception_body
               }
        \end{lstlisting}
        
        Exception termination constructs are allowed on several control structures
        - see \autorefph{sec:excep_term}.
        
        \subsection{immediate execution}
        
            On immediate execution this generates an in-line target statement.
            Immediate sequential execution is carried out through all the case
            expressions and the associated code bodies in order.
        
        \subsection{target execution}

            The branch test expression is evaluated and compared simultaneously
            with all case values (the case values are constant in the target
            code but may have been evaluated from expressions during immediate
            execution). If the branch test expression value matches a case value
            then the associated target statement is executed. The target
            statement may be a control structure containing nested target
            statements.

            If the test expression contains queue variables the test expression
            may wait on queue availability \index{queue!availability}. If the test
            expression reads a queue it will remove a datum. If a statement
            within an enclosed code body also reads the same queue, care must be
            taken to ensure that these reads occur simultaneously, otherwise two
            data will be read rather than one. In some cases it may be best to
            use an examine rather than a read on a queue variable within a test
            expression for a control structure if it is known that the queue
            variable is available for reading.

            The case selection takes no clock cycles but the branch test
            expression evaluation and parallel case comparison delay is added to
            the execution delay of every case block. If a case is not matched
            the {\em branch} statement will take one clock cycle).
            


    \section{the {\bf seq while}, {\bf par while} and {\bf sync while}
                                                        control structures}

        \index{statement!seq while}
        \index{statement!par while}
        \index{statement!sync while}
        This statement creates a {\bf while} loop in target code. The body
        of the loop is enclosed in a {\bf seq} block, a {\bf par} block or a
        {\bf sync} block.
        
        The format is -
        \begin{lstlisting}[gobble=8, language=threepl]
               seq while (expression) {
                   body
               }
        \end{lstlisting}
        or
        \begin{lstlisting}[gobble=8, language=threepl]
               par while (expression) {
                   body
               }
        \end{lstlisting}
        or
        \begin{lstlisting}[gobble=8, language=threepl]
               sync while (expression) {
                   body
               }
        \end{lstlisting}

        The expression must evaluate to a logical result. It is usually target
        mode but may be immediate mode. If the expression is immediate and
        false the body will not be subject to immediate execution and no
        target code will be generated for the statement. If the expression is
        immediate and true an infinite target loop will be created. Similarly
        the test expression may be a constant value mode variable in which
        case the behaviour is the same as for an immediate expression. If the
        expression is target mode and is not a constant value mode variable
        then it must have a single clock domain which is the same as the current
        clock domain.
        
        The code body is one or more statements.
        
        Any of the above forms of {\bf while} may be followed by an
        exception termination
        \begin{lstlisting}[gobble=8, language=threepl]
               unless (exception_expression) {}
        \end{lstlisting}
        or
        \begin{lstlisting}[gobble=8, language=threepl]
               unless (exception_expression) {
                   exception_body
               }
        \end{lstlisting}
        
        Exception termination constructs are allowed on several control structures
        - see \autorefph{sec:excep_term}.

        \subsection{immediate execution}
        
            On immediate execution this generates an in-line target statement.
            Immediate sequential execution is carried out on the code body
            of the {\em while} as long as the test expression is not immediate
            mode and false in value.
        
        \subsection{target execution}

            The {\bf seq while}, {\bf par while} and {\bf sync while} control
            structures repeatedly execute the code body while the test
            expression is true.

            When the expression and body are both available and the expression
            evaluates to true, control passes to the body.

            When the expression is available and evaluates to false the {\em
            while} will terminate.

            If the test expression contains queue variables the test expression may
            wait on queue availability \index{queue!availability}. If the test
            expression reads a queue it will remove a datum. If a statement within
            the enclosed code body also reads the same queue, care must be taken to
            ensure that these reads occur simultaneously, otherwise two data will
            be read rather than one. See the examples in the section on {\bf
            when()} (\autorefph{sec:when_struct} above) which also pertain to the
            target {\bf while()}.
            
            The compiler recognises cases where a target {\bf while()} contains
            queue reads or writes in the code body and also reads some but not
            all of these queues in the test expression, and in such cases issues
            a warning.
            
            The target {\em while} adds no extra cycles for each loop
            iteration test, but termination of the loop, even when there
            are no iterations, takes one extra clock cycle.

            Note that {\bf seq while}, {\bf par while} and {\bf sync while} control
            structures pass execution control through the test expression within the
            same cycle that the body executes. As a consequence the signal delay
            through the test expression may reduce the timing margins for the first
            statement (seq block) or statements (par block) in the code body. In
            some cases it may for example be possible to change the order of
            statements in a seq block so that a complex statement is not the first.

            If an exception termination construct is appended, execution within
            the {\bf while} can be prematurely terminated
            - see \autorefph{sec:excep_term}.



    \section{the {\bf seq do while}, {\bf par do while} and {\bf sync
                                            do while} control structures}

        \index{statement!seq do while}
        \index{statement!par do while}
        \index{statement!sync do while}
        This statement creates a {\bf do while} loop in target code. The body
        of the loop is enclosed in a {\bf seq} block, a {\bf par} block or a
        {\bf sync} block.

        The format is -
        \begin{lstlisting}[gobble=8, language=threepl]
               seq do {
                   body
               } while (expression);
        \end{lstlisting}
        or
        \begin{lstlisting}[gobble=8, language=threepl]
               par do {
                   body
               } while (expression);
        \end{lstlisting}
        or
        \begin{lstlisting}[gobble=8, language=threepl]
               sync do {
                   body
               } while (expression);
        \end{lstlisting}

        The expression must evaluate to a logical result and must be target
        mode.
        
        The code body is one or more statements.
        
        Any of the above forms of {\bf do while} may be followed by an
        exception termination
        \begin{lstlisting}[gobble=8, language=threepl]
               unless (exception_expression) {}
        \end{lstlisting}
        or
        \begin{lstlisting}[gobble=8, language=threepl]
               unless (exception_expression) {
                   exception_body
               }
        \end{lstlisting}
        
        Exception termination constructs are allowed on several control structures
        - see  \autorefph{sec:excep_term}.
        
        \subsection{immediate execution}
        
            On immediate execution this generates an in-line target statement.
            Immediate sequential execution is carried out on the code body
            of the {\em do while}.
        
        \subsection{target execution}

            The {\bf seq do while}, {\bf par do while} and {\bf sync do
            while} control structures repeatedly execute the body while the
            test expression remains true. This differs from the {\bf while}
            in that the body is executed at least once.
            
            If the code body contains queue reads and the test expression also
            reads the same queues care must be taken to ensure that where the
            code body and test expression are intended to read the same datum
            simultaneously the intent is not thwarted by unexpected waits in
            the body code. In particular if the test expression contains
            queue reads then the body code should either not read those queues
            or else should not contain any other queue reads or any queue writes.
            The test expression must have a single clock domain which is the
            same as the current clock domain.          
            
            Note that in the following -
            \begin{lstlisting}[gobble=12, language=threepl]
                   par do {
                       s <- p;
                   } while (p);
            \end{lstlisting}
            queue {\bf p} will be read on the initial execution of the body.
            In the next clock cycle the test expression will again read q1
            and if the result is true the body will execute on the same
            clock cycle reading queue {\bf p} simultaneously with the test
            expression.
                        
            The compiler recognises cases where a target {\bf do while()} contains
            queue reads or writes in the code body and also reads some but not
            all of these queues in the test expression, and in such cases issues
            a warning.
            
            The target {\em while} adds no extra cycles for each loop
            iteration test, but one extra clock cycle is required for
            loop termination.

            If an exception termination construct is appended, execution within
            the {\bf do while} can be prematurely terminated
            - see \autorefph{sec:excep_term}.



    \section{the {\bf seq\{\dots\}} sequential block}

        \index{statement!seq}
        The format is -
        \begin{lstlisting}[gobble=8, language=threepl]
               seq {
                   body
               }
        \end{lstlisting}
        The code body is one or more statements.
        
        The {\bf seq} may be followed by an exception termination
        \begin{lstlisting}[gobble=8, language=threepl]
               unless (exception_expression) {}
        \end{lstlisting}
        or
        \begin{lstlisting}[gobble=8, language=threepl]
               unless (exception_expression) {
                   exception_body
               }
        \end{lstlisting}
        
        Exception termination constructs are allowed on several control structures
        - see \autorefph{sec:excep_term}.
        
        \subsection{immediate execution}
        
            On immediate execution this generates an in-line target statement.
            Immediate sequential execution is carried out on the code body
            of the {\em seq}.
        
        \subsection{target execution}

            The {\bf seq} control structure executes the in-line target
            statements within it sequentially. The control structure
            finishes when the last contained statement finishes execution.
            
            The target {\em seq} adds no extra cycles to the time taken
            by the code body target statements.

            If an exception termination construct is appended, execution within
            the {\bf seq} can be prematurely terminated
            - see \autorefph{sec:excep_term}.



    \section{the {\bf par\{\dots\}} parallel block}

        \index{statement!par}
        The format is -
        \begin{lstlisting}[gobble=8, language=threepl]
               par {
                   body
               }
        \end{lstlisting}
        The code body is one or more statements.
        
        The {\bf par} may be followed by an exception termination

        \begin{lstlisting}[gobble=8, language=threepl]
               unless (exception_expression) {}
        \end{lstlisting}
        or
        \begin{lstlisting}[gobble=8, language=threepl]
               unless (exception_expression) {
                   exception_body
               }
        \end{lstlisting}
        
        Exception termination constructs are allowed on several control structures
        - see \autorefph{sec:excep_term}.
        
        \subsection{immediate execution}
        
            On immediate execution this generates an in-line target statement.
            Immediate sequential execution is carried out on the code body
            of the {\em par}.
        
        \subsection{target execution}

            The {\bf par} control structure executes the in-line target
            statements within it in parallel. The control structure
            finishes when the last contained statement finishes execution.
            
            The target {\em par} adds no extra cycles to the time taken
            by the longest code body target statement.

            If an exception termination construct is appended, execution within
            the {\bf par} can be prematurely terminated
            - see \autorefph{sec:excep_term}.



    \section{the {\bf sync\{\dots\}} synchronised parallel block}\label{sec:sync}

        \index{statement!sync}
        The format is -
        \begin{lstlisting}[gobble=8, language=threepl]
               sync {
                   body
               }
        \end{lstlisting}
        The code body is one or more statements.
        
        The {\bf sync} may be followed by an exception termination
        \begin{lstlisting}[gobble=8, language=threepl]
               unless (exception_expression) {}
        \end{lstlisting}
        or
        \begin{lstlisting}[gobble=8, language=threepl]
               unless (exception_expression) {
                   exception_body
               }
        \end{lstlisting}
        
        Exception termination constructs are allowed on several control structures
        - see \autorefph{sec:excep_term}.
        
        \subsection{immediate execution}
        
            On immediate execution this generates an in-line target statement.
            Immediate sequential execution is carried out on the code body
            of the {\em sync}.
        
        \subsection{target execution}

            The {\bf sync} control structure executes the in-line target statements
            within it in parallel but it will not begin parallel execution of
            the body until all queues read or written within the code body are
            available \index{queue!availability}. The control structure finishes
            when the last contained statement finishes execution.
            
            It is important to note that the {\bf sync} takes no account
            of conditional access to queues within its block, i.e. -
            \blist
            \item   queue reads or writes within a {\bf when} statement
            \item   queue reads within operands to the ternary conditional
                    operator \verb'?:'
            \blend
            If queue reads or writes are located within the block the {\bf sync}
            will wait for availability on all of these before execution of the
            target code in the block, even though some of these queue accesses
            may not eventuate due to conditional execution.
            
            The {\em sync} adds no extra cycles to the time taken
            by the code block unless it must wait for queue availability.
            
            Because the {\bf sync} block will not execute the code block
            within it until any queues referenced in that block are available,
            many statements which would otherwise test for queue availability
            before executing do not do so if they are within a {\bf sync}
            block. For example assignment statements immediately
            within a {\bf sync} block omit the queue availability test.
            Although the validity of execution is not affected by these
            considerations, using a {\bf sync}, rather than a {\bf par},
            to create a parallel block can result in savings in both
            logic components and path delays. The omission of
            queue availability tests within a {\bf sync} block does not extend -
            \blist
            \item   beyond the initial entry test of a {\bf seq while()\{\}}
                    or {\bf par while()\{\}} block
            \item   into a {\bf seq do while()\{\}} or {\bf par do while()\{\}}
                    block
            \item   beyond the first target statement of a {\bf seq} block
            \blend
            The omission of queue availability tests within a {\bf sync} block
            does extend into -
            \blist
            \item   when blocks
            \item   par and sync blocks
            \item   for, while and do while blocks
            \blend
            Generally it is good practice to use a {\bf sync} instead of
            a {\bf par} when queues are used and it is not necessary to
            wait on individual availability of those queues. It is also
            preferable to enclose a {\bf when} in a {\bf sync} rather than
            use {\bf sync} blocks within the {\bf when}, e.g.
            \begin{lstlisting}[gobble=12, language=threepl]
                   sync {
                       when (...) par {
                           a <<- c;
                           b <<- d;
                       } else par {
                           a <<- e;
                           b <<- f;
                       }
                   }
            \end{lstlisting}
            or
            \begin{lstlisting}[gobble=12, language=threepl]
                   sync {
                       when (...) sync {
                           a <<- c;
                           b <<- d;
                       } else sync {
                           a <<- e;
                           b <<- f;
                       }
                   }
            \end{lstlisting}
            are preferable to
            \begin{lstlisting}[gobble=12, language=threepl]
                   when (...) sync {
                       a <<- c;
                       b <<- d;
                   } else sync {
                       a <<- e;
                       b <<- f;
                   }
            \end{lstlisting}
            In the last case each assignment statement will contain logic to
            test for the availability of the queues concerned, whereas the
            previous two cases will test only once for all queues before
            executing the {\bf when}. However, note that this might not be
            desirable where the alternative blocks in the {\bf when} handle
            different queues. In the last case the {\bf when} might select the
            true block and execute it when queues {\em a}, {\em b}, {\em c} and
            {\em d} are available regardless of the availability of queues {\em
            e} and {\em f}. The other two cases would not enter the outer sync
            unless all six queues were available even though two of them would
            not be used.

            If an exception termination construct is appended, execution within
            the {\bf sync} can be prematurely terminated
            - see \autorefph{sec:excep_term}.

    \section{the {\bf [\dots]} inline module}

        \index{module!in-line}
        The format is -
        \begin{lstlisting}[gobble=8, language=threepl]
               [
                   body
               ]
        \end{lstlisting}
        The code body is one or more statements.
        
        \subsection{immediate execution}
        
            On immediate execution this generates an inline module.
            Immediate sequential execution is carried out on the code body.
        
        \subsection{target execution}

            The inline module created is separate from the module in
            which it is created though it shares the scope of the
            module, procedure or function in which it is created. It does not
            generate any in-line target code in the module, procedure or
            function in which it is created.


    \section{the exception termination construct}\label{sec:excep_term}
        
        \index{target!exception termination}
        \index{unless construct}
        \index{unless construct}
        The control structures -
        \blist
        \item   when
        \item   branch
        \item   seq
        \item   par
        \item   sync
        \item   sync when
        \item   seq while
        \item   par while
        \item   sync while
        \item   seq do while
        \item   par do while
        \item   sync do while
        \blend        
        may be followed by an exception termination
        \begin{lstlisting}[gobble=8, language=threepl]
               unless (exception_expression) {}
        \end{lstlisting}
        or
        \begin{lstlisting}[gobble=8, language=threepl]
               unless (exception_expression) {
                   exception_body
               }
        \end{lstlisting}
        
        The exception body is a single target statement but this might be a
        {\bf seq}, {\bf par} or other type of target code block statement.
                
        A control structure with an exception termination cannot be nested
        within another control structure with an exception termination, either
        in a control structure code block or an exception code block\footnote{
        This restriction may be removed in the future.}.
        
        \subsection{immediate execution}
        
            Following immediate execution of the code within the control structure
            the code in the exception body, if present, is executed. An exception
            body cannot execute a {\bf break}, {\bf continue} or {\bf return} (but
            a {\bf break} or {\bf continue} could be nested within a {\bf for} loop
            in the exception body for example).
        
        \subsection{target execution}
        
            When the exception expression is true, any target code executing
            within the associated control structure will cease, including
            execution within any nested target control structures. The clock
            cycle in which the exception expression is true will override any
            statements executing in the control structure. If there is an
            exception code block, that code will then be executed. The control
            structure then terminates as though it had executed to completion
            normally.

            An exception termination takes a minimum of one cycle.

            Note that if code within a control structure has executed a {\bf
            stop()} it would never ordinarily complete. An exception termination
            will force completion of that control structure regardless of the
            stop(), i.e. the code block will skip to the end despite any
            intervening stop() even if the code has already reached the stop().

            An exception block may contain a {\bf stop()} in which case
            execution of the exception block will cease at that point and the
            associated control structure will not complete.

            If the exception expression does not contain queue reads, it may
            remain true for an indefinite time. Execution of the associated
            control structure will cease immediately and any exception code will
            be executed. Completion of the associated control structure will
            then occur if the exception expression is now false, otherwise
            completion will be delayed until the exception expression becomes
            false. By prolonging the exception expression true value completion
            of the associated control structure can be delayed. Any excursions
            of the exception expression value during execution of the exception
            block will be ignored.

            If the exception expression contains queue reads then the exception
            condition will occur when a queue read occurs and the resulting value
            is true. False values are read and ignored. Execution of the
            associated control structure will cease immediately, any exception
            code will be executed and completion of the associated control
            structure will then occur. If the exception expression is available
            on completion the expression will again be evaluated and if true
            another exception will occur.

            An exception code block may contain queue reads and writes. These
            will wait on queue availability as usual.

        \subsection{notes}
        
            Cessation of execution on an exception termination includes any
            statements waiting on queues or {\bf waitpri} statements waiting on
            priority. The queue variables themselves are not affected, a pending
            read or write simply failing to execute. A single assignment
            involving queues can be terminated on an exception by enclosing it
            in a {\bf seq}, e.g.
            \begin{lstlisting}[gobble=12, language=threepl]
                   seq {
                       q1 <<- q2 + 2;
                   } unless (reset) {}
            \end{lstlisting}
            but this example is pointless as the statement would simply
            continue with another execution after the reset. A more likely
            scenario would be to clear a processing pipeline when the exception
            occurs, e.g.
            \begin{lstlisting}[gobble=12, language=threepl]
                   seq {
                       q1 <<- q2 + 2;
                   } unless (reset) {
                       par {
                           reset(q1, q2);
                           .
                           . other resetting code
                           .
                       }
                    }
            \end{lstlisting}
            The pending assignment will be killed, the queues cleared and other
            reset code executed, then, if at the top level, the {\bf seq} will
            resume execution.
            This might alternatively have been written
            \begin{lstlisting}[gobble=12, language=threepl]
                   seq {
                       q1 <<- q2 + 2;
                   } unless (reset) {}
                   
                   when (reset) par {
                       reset(q1, q2);
                       .
                       . other resetting code
                       .
                   }
            \end{lstlisting}
            Note that in this case if {\bf reset} is true for one cycle the {\bf
            seq} will execute again on the following cycle if it is at the top
            level, since there is no exception code block to delay completion,
            but the queues will have been cleared on the same clock cycle as the
            reset prior to continued execution of the {\bf seq}.

    \section{the waitpri () construct}\label{sec:waitpri}
    
        \index{waitpri}
        \index{wait on priority}
        The {\bf waitpri} control structure allows arbitration between
        conflicting statements, usually but not necessarily involving queue
        writes or memory reads or writes. It takes the forms -
        \begin{lstlisting}[gobble=8, language=threepl]
               waitpri (pv, p)
                   statement
        \end{lstlisting}
        or
        \begin{lstlisting}[gobble=8, language=threepl]
               waitprilock (pv, p)
                   statement
        \end{lstlisting}
        where {\bf pv} is a priority mode variable and {\bf p} is an immediate
        mode expression with an integer value. {\bf p} can have a non-zero
        positive value where lower values have a higher priority (zero is
        reserved). {\bf p} has a maximum value which is one less than the
        dimension of the priority variable array (default dimension is 20).
        Priority values for a particular priority variable do not have to be
        contiguous but must be unique.
        
        The priority variable must not be assigned anywhere (the {\bf waitpri
        ()} and {\bf waitprilock} perform assignments to the priority variable).
        The values of the priority variable boolean array may be used elsewhere,
        each member being true when the associated statement executes. The first
        entry in the priority array, subscript 0, is used to lock out priority
        requests when execution has been granted by {\bf waitprilock}. This
        output will go high on the clock cycle following execution granting at
        any priority level and will remain high until inbuilt target procedure
        {\bf priunlock()} is called.
        
        In the case of {\bf waitprilock} if there are no enclosed target
        statements (or perhaps no statements at all, i.e. just a ";") then a
        single cycle {\bf nop()} is provided. The {\bf waitpri} must have a
        single target statement (which might be a {\bf sync} statement or a {\bf
        sync when} statement containing multiple target statements).
        
        A {\bf waitpri} requests granting of execution of its enclosed
        statement via the priority variable and the specified priority
        level. Where queue availability is required this is waited upon
        if necessary prior to requesting priority. Priority is only
        held for one clock cycle.
        
        The {\bf waitprilock} is the same as {\bf waitpri} but it continues
        to hold priority. Priority is released by executing the inbuilt
        inline target code procedure {\bf priunlock()}
        \autorefph{sec:ippriunlock}. {\bf priunlock()} takes one clock cycle to
        execute but can usually be executed in parallel with the last
        statement to be executed under mutual exclusion.
        
        The same priority variable may be used by both {\bf waitpri} and
        {\bf waitprilock} statements.
        
        By default each priority level for a priority variable may be
        used only once. A situation may occur where it is desired to have
        two or more waitpri or waitprilock statements waiting on the same
        priority where it is known that the execution logic ensures that
        they are mutually exclusive. In that case the default can be
        overridden by setting the attribute "multiplein" true on the
        priority variable.

        The {\bf waitpri} or {\bf waitprilock} will delay execution until a
        clock cycle on which -
        \blist
        \item   all queue variables in the statement are available
        \item   the priority variable grants execution, i.e. no higher
                priority request has pre-empted it
        \blend
        These conditions are not latched; queue availability may come and
        go cycle by cycle as may priority pre-emption. If these conditions
        are met immediately then execution is not delayed, i.e. no clock cycle
        overheads are imposed unless queue non-availability or priority
        pre-emption intervenes.

        The statement enclosed within a {\bf waitpri} or {\bf waitprilock}
        may only be -
        \blist
        \item   a target assignment statement
        \item   a {\bf sync} block
        \item   a {\bf par} block which contains no queue reads or writes
        \item   cmemorywrite()
        \item   rmemoryread()
        \item   rmemorywrite()
        \item   nop() (restricted to single cycle)
        \item   assert()
        \item   reset()
        \blend
        
        Inbuilt inline target procedure {\bf stop()} is not allowed.
        
        Nested sync or par blocks in turn must also only contain
        statements from the same restricted list.

        The {\bf waitpri} only guarantees priority for a single cycle so
        execution must occur in one cycle. Thus the enclosed statement must be
        controlled at a single point. At this point in execution a request is
        made for priority to be granted. Thus a single assignment is OK and a
        {\bf sync} block also controls execution of an enclosed block at a
        single point (where it waits on queue availability). A {\bf par} block
        will also exxecute in a single clock cycle as long as it does not read
        or write a queue mode variable.

        These restrictions are enforced by the compiler.
        
        {\bf Note that for a waitprilock these restrictions only apply to the
        statement associated with the waitprilock, not to following
        statements up to the priunlock() call!}
        
        A {\bf waitpri} or {\bf waitprilock} should not be nested within -
        
        \blist
        \item   a {\bf sync} block
        \item   a {\bf sync when } block
        \blend
        
        however it is not an error. A sync block is intended to ensure that its
        enclosed statements (or at least the first of a nested {\bf seq}) can
        execute simultaneously. An enclosed {\bf waitpri} or {\bf waitprilock}
        might not execute because of priority pre-emption. While this would not
        be an error it rather defeats the purpose of the {\bf sync} since the
        delay means that queue availability would again have to be tested.
        
        A {\bf waitpri} or {\bf waitprilock} cannot be nested within another
        {\bf waitpri} or {\bf waitprilock} as the statement enclosed in the
        outer one must be able to execute when enabled but the inner one may
        preclude that. Again the compiler prohibits this.
        
        Where a {\bf sync} is nested within a {\bf waitpri} or {\bf
        waitprilock}  the {\bf waitpri} or {\bf waitprilock} actually passes
        control of priority seeking to the enclosed {\bf sync} so that the sync
        waits on both queue availability (if there are queue reads or writes in
        its body) and priority granting.
        
        The following example arbitrates writes to a queue -
        \begin{lstlisting}[gobble=8, language=threepl]
               priority(pv);
               queue({p, q1, q2}, t);
               
               waitpri (pv, 1)
                   p <<- q1;
               waitpri (pv, 2)
                   p <<- q2;
        \end{lstlisting}
        Data appearing in queues {\bf q1} or {\bf q2} will be written to
        queue {\bf p}. Where they arrive simultaneously {\bf q1} will prevail.
        If multiple assignments are required then a {\bf sync} block may be used -
        \begin{lstlisting}[gobble=8, language=threepl]
               priority(pv);
               queue(p, "[2]t");
               queue({q1, q2, q3}, t);
               static(s, t);

               waitpri (pv, 1)
                   sync {
                       p[0] <<- q1;
                       p[1] <<- q2;
                   }
               waitpri (pv, 2)
                   sync {
                       p[0] <<- q3;
                       p[1] <<- s;
                   }
        \end{lstlisting}
        If sequential statements are to be executed while holding
        priority then a waitprilock must be used -
        \begin{lstlisting}[gobble=8, language=threepl]
               priority(pv);
               queue(p, t);
               queue({q1, q2, q3}, t);
               static(s, t);

               seq {
                   waitprilock (pv, 1)
                       p <<- q1;
                   p <<- q2;
                   sync {
                       p <<- q3;
                       priunlock(pv);
                   }
               }
        \end{lstlisting}
        
        A {\bf waitprilock} may be used as an execution gate to inhibit
        processing of queue data until downstream processing of the previous
        datum has finished. In that case the priunlock() may be executed
        in a different execution thread -
        \begin{lstlisting}[gobble=8, language=threepl]
               priority(pv);
               queue(p, "[4]uint:8");
               queue(q1, "log");

               seq {
                   static(x, "[4]uint:8");
                   static({count, addr, i}, "uint:3");
                   
                   waitprilock (pv, 1)
                       sync {
                           x <- p;
                           count <- 4;
                           addr <- 0;
                           i <- 0;
                       }
                   par while (count != 0) par {
                       rmemorywrite(m, 0, addr, x[i]);
                       addr++;
                       count--;
                       i++;
                   }
                   q1 <<- true;
               }
               
               seq {
                   when (q1)
                   .
                   .
                   .
                   rmemoryread(m, 1, 0);
                   par while (...) {
                       ... <- r memoryout(m, 1);
                       rmemoryread(m, 1, a);
                   }
                   .
                   .
                   priunlock(pv);
               }
        \end{lstlisting}
        The first {\bf seq} will attempt to repeatedly execute as long as
        queues {\bf p} and {\bf q1} are available. The second {\bf seq}
        repeats, waiting on queue {\bf q1} each time. The assumption here is
        that the first {\bf seq} alters state that is shared with the second
        {\bf seq} and this requires that data be received singly. In
        this example each repeat of the first {\bf seq} will wait at the {\bf
        waitprilock()} until the second {\bf seq} executes the {\bf
        priunlock()} at the end of processing the previous datum.

        If the data sources are not queues, conditional statements may
        be used -
        \begin{lstlisting}[gobble=8, language=threepl]
               priority(pv);
               queue(p, t);
               static({s1, s2}, t);
               input(l1, "log", ........);
               input(l2, "log", ........);
               
               when (l1)
                   waitpri (pv, 1)
                       p <<- s1;
               when (l2)
                   waitpri (pv, 2)
                       p <<- s2;
        \end{lstlisting}
        Note that a {\bf when} cannot be nested within a {\bf waitpri}.
        
        {\bf Waitpri} is also useful for arbitration among memory accesses
        where the arguments to the inbuilt procedure call may or may not
        involve queues -
        \begin{lstlisting}[gobble=8, language=threepl]
               priority(pv);
               rmemory(m, 1, atype, dtype);
               static(addr1, atype);
               queue(addr2, atype);
               static(data, dtype);
               value(l, "log");
               
               when (l)
                   waitpri (pv, 1)
                       rmemorywrite(m, 0, addr1, data);
               waitpri (pv, 2)
                   rmemoryread(m, 0, addr2);
        \end{lstlisting}
        
        {\bf Waitpri} may also be used for mutually exclusive control of
        a bus -
        \begin{lstlisting}[gobble=8, language=threepl]
               priority(pv);
               selectvalue(bus, t);
               static(s, t);
               value(l, "log");
               queue(p, t);
               
               when (l)
                   waitpri (pv, 1)
                       bus <<- s;
               waitpri (pv, 2)
                   bus <<- p;
        \end{lstlisting}
        Here when {\bf l} is true static variable {\bf s} will be written to the
        bus for one clock cycle. In this case no queue waits are involved. When
        queue {\bf p} is available it will be written to the bus when not
        pre-empted by the other write.
        
        

    \section{the petrinet construct}
    
        \index{petrinet}
        A Petri net contains a number of nodes known as {\em places}. Each
        place can hold one token (a {\em finite capacity} Petri net where each
        place has a capacity of 1). A single token can be taken from
        one or more source places and a single token put in one or
        more destination places as a result of a transition
        (an {\em ordinary} petri net). A transition can also additionally
        depend on a logical {\em predicate}. The Petri net must be
        {\em marked} to initially place tokens in some places.
        
        A Petri net where only one place is marked and each transition
        has one source place and one destination place is a simple state
        machine.
        
        The Petri net is effectively a module in that the generated logic is
        not in-line and the Petri net acts as a single queue destination for
        queue acknowledgement.
        
        The syntax is illustrated by the following example -
        \begin{lstlisting}[gobble=8, language=threepl]
               petrinet (r) {
                   initial q1, q2;
                   q1, q2 -> q3, q4 when (a && (b != 0)) {i = 1; s <- queue1;};
                   q3 ->> q1 {l1 <- queue1;};
                   q4 -> q2, q3 when (!l2);
               }
        \end{lstlisting}
        Here the Petri net has four places, {\bf q1} to {\bf q4}. If these
        identifiers have already been created they must be mode static
        and type log and must not have been previously assigned. If they
        have not been created {\bf petrinet} will create them. Normal scope
        rules apply to these variables. Place variables cannot be assigned
        values except via the Petri net mechanism. The values of place
        variables can be used as for normal variables.
        
        The first line contains \verb'(r)' which is an optional (parenthesised)
        reset expression. If present a logical true value of this expression
        will reset the Petri net to its initial marking. This will override any
        concurrent transitions. The reset will not affect any transition code
        blocks that are currently executing.
        
        The second line must be present and lists the place marking. In this
        case places {\bf q1} and {\bf q2} will initially contain tokens (will
        be true).
        
        The following lines describe allowed transitions. A transition
        is a list of initial (source) places, a right single or double arrow,
        a list of final (destination) places, an optional predicate and
        finally an optional code block to be executed on transition. Each
        transition statement is terminated by a semicolon.
        
        A single right arrow specifies a {\em weak} transition. The transition
        occurs when the source places each contain a token (are true) and the
        predicate, if present, is true. The transition results in the source
        places losing their tokens (all become false) and the destination places
        each receiving a token (all become true). With a weak transition a
        destination place may already have a token but since it can only hold
        one token it still contains a single token after the transition.
              
        A double right arrow specifies a {\em strict} transition. This
        has the additional constraint over a weak transition that the
        transition can only occur if none of the destination places have
        a token.
        
        When a transition occurs an associated code block will be
        simultaneously executed. If the code block takes longer than one
        clock cycle to execute then the transition associated with the code
        block cannot occur again until the code block has finished execution.
        This does not inhibit any other transitions, including transitions
        which affect the same places.
        
        The optional predicate is an additional condition which must be met
        for a transition to fire. It may contain any logical expression and that
        expression may include place variables. Note that all place
        variables are created prior to analysing the transitions hence
        they can be used in transition predicates ahead of any other occurrences.

        A transition may be duplicated if the predicates differ. This
        effectively ORs the predicates but allows each transition to
        execute a separate code block.

        Predicates may contain queue reads but these must be used with
        caution. The {\bf when} clause specifying the predicate executes
        every clock cycle and thus will execute any queue reads whenever
        those queues are available regardless of whether a transition is
        possible. It only makes sense to use a queue read in a predicate
        if the transition is always ready to fire when the queue becomes
        available.
        
        Transition execution blocks may also contain queue reads (or writes).
        The Petri net as a whole is a module and so queue reads wherever they
        occur within the Petri net will be acknowledged individually. The
        Petri net will act as a single queue destination and a queue source
        will wait until a read has occurred within the Petri net before
        advancing to the next datum. Queue reads occurring anywhere within
        the Petri net, either in predicates or code blocks, will constitute
        a single read if they occur simultaneously or as separate reads
        otherwise.
        
        A transition code block may contain any number of immediate statements but
        must contain exactly one in-line target statement. The latter of course may
        be a {\bf seq}, {\bf par} etc. which nest other target statements.
        Immediate statements within code blocks are executed sequentially from top
        to bottom of the Petri net construct. Care must be taken if a transition
        code block executes for more than two cycles. A transition cannot occur if
        an associated transition code block is executing (because of a previous
        occurence of the same transition).




 
